3. Pruebe que los siguientes números no son racionales. Consejo: suponga que sí en cada caso. \\

\textit{Definiremos que los racionales son números de la forma $\frac{p}{q} : p,q\in \mathbb{N}$} con la condicion $\gcd(p, q) = 1$, i.e., primos relativos o fracciones irreducibles.

{a) } \(\sqrt{3}\) \text{ no es racional.}

\text{Supongamos que } \(\sqrt{3} = \frac{p}{q}\) \text{ con } \(p, q \in \mathbb{Z}, \gcd(p,q) = 1\).

\[
3q^2 = p^2.
\]

\text{Entonces } \(3 \mid p^2 \implies 3 \mid p\).

\text{Sea } \(p = 3k\), \text{ entonces }

\[
3q^2 = 9k^2 \implies q^2 = 3k^2 \implies 3 \mid q.
\]

\text{Esto contradice que } \(\gcd(p,q) = 1\).

\medskip

{b) } \(\sqrt{5}\) \text{ no es racional.}

\text{Procedemos igual: supongamos } \(\sqrt{5} = \frac{p}{q}\) \text{ con } \(\gcd(p,q) = 1\).

\[
5q^2 = p^2 \implies 5 \mid p.
\]

\text{Sea } \(p = 5k\), \text{ entonces }

\[
5q^2 = 25k^2 \implies q^2 = 5k^2 \implies 5 \mid q,
\]

\text{contradicción.}

\medskip

{c) } \(\sqrt{24}\) \text{ no es racional.}

\text{Notemos que } \(\sqrt{24} = 2\sqrt{6}\).

\text{Si } \(\sqrt{24}\) \text{ fuera racional, entonces } \(\sqrt{6}\) \text{ también lo sería.}

\text{Supongamos } \(\sqrt{6} = \frac{p}{q}, \gcd(p,q) = 1\).

\[
6q^2 = p^2.
\]

\text{Entonces } \(2 \mid p \text{ y } 3 \mid p \implies 6 \mid p\).

\text{Sea } \(p = 6k\), \implies 

\[
q^2 = 6k^2 \implies 6 \mid q,
\]

\text{contradicción.}

\medskip

{d) } \(2^{1/3}\) \text{ no es racional.}

\text{Supongamos } \(2^{1/3} = \frac{p}{q}, \gcd(p,q) = 1\).

\[
2 q^3 = p^3.
\]

\text{Entonces } \(2 \mid p\).

\text{Sea } \(p = 2k\), \implies 

\[
2 q^3 = 8 k^3 \implies q^3 = 4 k^3 \implies 2 \mid q,
\]

\text{contradicción.}

\medskip

{e) } \(13^{1/4}\) \text{ no es racional.}

\text{Supongamos } \(13^{1/4} = \frac{p}{q}, \gcd(p,q) = 1\).

\[
13 q^4 = p^4.
\]

\text{Entonces } \(13 \mid p\).

\text{Sea } \(p = 13k\), \implies 

\[
13 q^4 = 13^4 k^4 \implies q^4 = 13^3 k^4 \implies 13 \mid q,
\]

\text{contradicción.}

\medskip

{f) } \((5 - \sqrt{3})^{1/3}\) \text{ no es racional.}

\text{Supongamos que } \(x = (5 - \sqrt{3})^{1/3} \in \mathbb{Q}\).

\[
x^3 = 5 - \sqrt{3}.
\]

\text{Pero } \(5 - \sqrt{3} \notin \mathbb{Q}\), \text{ por lo que } \(x \notin \mathbb{Q}\).

\medskip

{g) } \((2 + \sqrt{2})^{1/2}\) \text{ no es racional.}

\text{Supongamos } \(y = (2 + \sqrt{2})^{1/2} \in \mathbb{Q}\).

\[
y^2 = 2 + \sqrt{2}.
\]

\text{De aquí } \(\sqrt{2} = y^2 - 2 \in \mathbb{Q},
\)

\text{lo cual es falso.}

\medskip

\text{Por lo tanto, ninguno de los números dados es racional.} \\\\
